%==============================================================================
\section{DBSCAN Clustering}
\label{sec:dbscan}
%==============================================================================

\begin{dinhnghia}[title={DBSCAN}]
\textbf{D}ensity-\textbf{B}ased \textbf{S}patial \textbf{C}lustering of \textbf{A}pplications with \textbf{N}oise:
thuật toán phân cụm dựa trên mật độ, cần hai tham số --- \texttt{minPts} (số láng giềng tối thiểu) và \texttt{eps} (khoảng cách tối đa giữa điểm lõi).
\end{dinhnghia}

\subsection{Cách DBSCAN hoạt động}

\begin{phuongphap}[title={Các bước thuật toán}]
\begin{enumerate}
  \item Chọn ngẫu nhiên một điểm dữ liệu chưa được thăm.
  \item Nếu điểm có ít nhất \texttt{minPts} láng giềng trong bán kính \texttt{eps}: tạo cụm mới, thêm điểm và các láng giềng vào cụm.
  \item Nếu không đủ láng giềng: đánh dấu điểm là \textbf{nhiễu} (noise) và chuyển sang điểm tiếp theo.
  \item Lặp các bước 1--3 cho đến khi mọi điểm đã được thăm.
\end{enumerate}
\end{phuongphap}

\subsection{Triển khai Python (scikit-learn)}

\begin{vidu}[title={Bước 1: Nạp dataset make\_moons}]
\begin{lstlisting}
from sklearn.datasets import make_moons
X, y = make_moons(n_samples=200, noise=0.05, random_state=0)
\end{lstlisting}
\end{vidu}

\begin{vidu}[title={Bước 2: DBSCAN với eps và min\_samples}]
\begin{lstlisting}
from sklearn.cluster import DBSCAN
clustering = DBSCAN(eps=0.2, min_samples=5)
clustering.fit(X)
\end{lstlisting}
\end{vidu}

\begin{vidu}[title={Bước 3: Trực quan hóa kết quả}]
\begin{lstlisting}
import matplotlib.pyplot as plt
plt.figure(figsize=(7.5, 3.5))
plt.scatter(X[:, 0], X[:, 1], c=clustering.labels_, cmap='rainbow')
plt.show()
\end{lstlisting}
\end{vidu}

\begin{vidu}[title={Ví dụ đầy đủ}]
\begin{lstlisting}
from sklearn.datasets import make_moons
from sklearn.cluster import DBSCAN
import matplotlib.pyplot as plt

X, y = make_moons(n_samples=200, noise=0.05, random_state=0)
clustering = DBSCAN(eps=0.2, min_samples=5)
clustering.fit(X)

plt.figure(figsize=(7.5, 3.5))
plt.scatter(X[:, 0], X[:, 1], c=clustering.labels_, cmap='rainbow')
plt.show()
\end{lstlisting}
\end{vidu}

\begin{ynghia}[title={Kết quả mong đợi}]
Biểu đồ phân tán cho thấy hai cụm tách bạch (hai ``vầng trăng''); điểm nhiễu thường được tô màu đen (nhãn $-1$).
\end{ynghia}

\tsimg{08_dbscan/img01.jpg}

\subsection{Ưu và nhược điểm}

\begin{tinhchat}[title={Ưu điểm}]
\begin{itemize}
  \item Xử lý cụm \textbf{hình dạng tùy ý}, không giả định cụm hình cầu như K-means.
  \item \textbf{Không cần biết trước} số cụm.
  \item \textbf{Phát hiện outlier}: điểm xa vùng mật độ cao được coi là nhiễu.
  \item Tương đối \textbf{ít nhạy} với khởi tạo so với K-means.
  \item \textbf{Mở rộng} tốt: chỉ cần khoảng cách giữa các láng giềng, không phải mọi cặp điểm.
\end{itemize}
\end{tinhchat}

\begin{luuy}[title={Nhược điểm}]
\begin{itemize}
  \item \textbf{Nhạy} với \texttt{eps} và \texttt{min\_samples}; chọn sai có thể gộp cụm hoặc bỏ sót cụm.
  \item Kém trên dataset \textbf{mật độ không đồng nhất} (giả định mật độ cụm tương tự).
  \item Có thể \textbf{không xác định} giữa các lần chạy.
  \item \textbf{Tốn kém} trên dữ liệu chiều cao (tính khoảng cách đắt).
  \item Nhiễu/outlier mật độ cao có thể bị gán nhầm vào cụm.
\end{itemize}
\end{luuy}
